% chapters/screenshot.tex
\chapter{Tangkapan Layar Dashboard}

Berikut adalah tangkapan layar dari bagian-bagian utama dashboard \textit{``Rokok Nomor Satu, Gizi Lain Waktu''}.

\section{Tampilan Utama -- Halaman Indonesia}

\begin{figure}[H]
  \centering
  \includegraphics[width=0.97\textwidth]{images/home_full}
  \caption{Halaman Indonesia -- tampilan utama dashboard dengan KPI strip, peta choropleth 38 provinsi, donut piring pengeluaran, dan kontrol global (dropdown metrik dan region)}
\end{figure}

Halaman Indonesia menampilkan empat komponen utama:
\begin{itemize}
  \item \textbf{KPI Strip}: 5 kartu metrik kunci -- jumlah provinsi aktif (34/38), rata-rata rasio rokok/gizi (38.7\%), rata-rata perokok harian (20.4\%), rata-rata stunting balita (16.1\%), dan provinsi dengan rasio tertinggi (Sulawesi Barat).
  \item \textbf{Peta Choropleth}: Visualisasi 38 provinsi Indonesia diwarnai berdasarkan rasio rokok terhadap total pengeluaran gizi. Warna merah tua menandai rasio tertinggi; 4 provinsi pemekaran Papua ditampilkan abu-abu (\textit{greyed-out}) karena data pengeluaran belum tersedia.
  \item \textbf{Donut Piring Pengeluaran}: Proporsi nasional pengeluaran rokok (28\%) dibandingkan 5 komponen gizi -- setiap komponen dikodekan dengan warna berbeda agar dapat dibedakan.
  \item \textbf{Kontrol Global}: Dropdown metrik (4 pilihan) dan dropdown region (7 wilayah) memperbarui seluruh visualisasi secara sinkron.
\end{itemize}

\section{Ranking Bar dan Butterfly Chart}

\begin{figure}[H]
  \centering
  \includegraphics[width=0.97\textwidth]{images/home_butterfly}
  \caption{Ranking bar 10 provinsi tertinggi (atas) dan butterfly chart karakteristik sosial-ekonomi (bawah) -- dimensi pendidikan ditampilkan; dropdown di kanan memungkinkan pengguna mengganti dimensi}
\end{figure}

Dua visualisasi tambahan di halaman Indonesia:
\begin{itemize}
  \item \textbf{Ranking Bar}: 10 provinsi dengan rasio rokok/gizi tertinggi. Garis putus-putus emas menandai rata-rata nasional. Panel \textit{Insight} di kanan menampilkan narasi dinamis sesuai filter aktif.
  \item \textbf{Butterfly Chart}: Membandingkan prevalensi perokok harian (kiri, merah) dengan persentase gizi normal balita (kanan, hijau) per kelompok karakteristik. Pengguna dapat mengganti dimensi melalui dropdown (Jenis Kelamin, Pendidikan, Pekerjaan, Tempat Tinggal, Status Ekonomi).
\end{itemize}

\section{Halaman Provinsi -- Profil dan Diagnosa}

\begin{figure}[H]
  \centering
  \includegraphics[width=0.97\textwidth]{images/province_top}
  \caption{Halaman Provinsi (Nusa Tenggara Timur) -- badge risiko, dua baris KPI, leaderboard struktur pengeluaran, dan kompas kuadran posisi provinsi}
\end{figure}

Halaman provinsi menampilkan:
\begin{itemize}
  \item \textbf{Header}: Nama provinsi, badge risiko berwarna (Risiko Tinggi/Sedang/Rendah berdasarkan rasio rokok/gizi), dan label region.
  \item \textbf{KPI Baris 1}: Rokok/kapita, Gizi total, Rasio Rokok \% Gizi, Protein/kapita -- dikodekan warna merah (rokok), hijau (gizi), emas (rasio), dan putih.
  \item \textbf{KPI Baris 2}: Perokok harian, Stunting balita, Mantan perokok, Rank nasional.
  \item \textbf{Leaderboard Bar}: Struktur 6 komponen pengeluaran (rokok + 5 gizi) untuk provinsi terpilih -- rokok ditampilkan dengan opasitas penuh sebagai penanda.
  \item \textbf{Kompas Kuadran}: Scatter plot posisi provinsi terhadap rata-rata nasional (sumbu X: rokok/gizi, sumbu Y: protein/kapita). Empat zona berwarna membantu pembacaan risiko: pojok kanan bawah = tekanan ganda.
\end{itemize}

\section{Simulator Kebijakan}

\begin{figure}[H]
  \centering
  \includegraphics[width=0.97\textwidth]{images/province_simulator}
  \caption{Simulator kebijakan -- slider realokasi pengeluaran rokok, tiga kartu estimasi (stunting, dana dihemat, setara konversi), dan Sankey diagram aliran dana dengan slider strategi alokasi vertikal}
\end{figure}

Simulator kebijakan merupakan fitur interaktif unggulan dashboard ini. Pengguna dapat:
\begin{itemize}
  \item Menggeser slider untuk mengatur jumlah pengeluaran rokok (0 hingga 2,5$\times$ nilai saat ini).
  \item Melihat estimasi perubahan stunting balita berdasarkan model regresi linear lintas 34 provinsi (R\textsuperscript{2}=0.21).
  \item Melihat besaran dana yang ``dihemat'' dan setara konversinya (butir telur, porsi ikan).
  \item Mengatur strategi alokasi dana hemat melalui slider vertikal (Komposisi / Rata / Protein / Serat), yang mengubah distribusi Sankey diagram secara \textit{real-time}.
\end{itemize}

\textit{Catatan: Semua estimasi bersifat asosiasi statistik lintas provinsi, bukan prediksi kausal individual. Disclaimer ini ditampilkan eksplisit di bawah Sankey.}

\section{Fitur Interaktif Utama}

\begin{table}[H]
\centering
\caption{Ringkasan fitur interaktif dashboard}
\renewcommand{\arraystretch}{1.3}
\begin{tabularx}{\textwidth}{@{}lXl@{}}
\toprule
\textbf{Fitur} & \textbf{Deskripsi} & \textbf{Halaman} \\
\midrule
Dropdown Metrik & Mengubah metrik peta (4 pilihan) & Indonesia \\
Dropdown Region & Menyaring data per wilayah (7 pilihan) & Indonesia \\
Klik Peta & Navigasi ke profil provinsi yang diklik & Indonesia \\
Hover Tooltip & Menampilkan detail provinsi saat kursor diarahkan & Indonesia \\
Dropdown Butterfly & Mengganti dimensi karakteristik (5 pilihan) & Indonesia \\
Slider Rokok & Mengatur skenario realokasi pengeluaran & Provinsi \\
Slider Alokasi & Mengubah strategi pembagian dana (4 mode) & Provinsi \\
Zoom \& Pan Peta & Zoom in/out dan geser peta interaktif & Indonesia \\
\bottomrule
\end{tabularx}
\end{table}
